\section{Introduction}

% A \textit{language model} assigns probabilities of sequences of words from the vocabulary $V$; the number of distinct sentences increases exponentially w.r.t to their length $N$. Hence the domain of a language model is, by definition, the exponential space $V^N$.  However, due to the vanilla exponential space being intractable, existing work tends to use recurrent architectures to generate conditional probabilities based on the context (typically encapsulated as a fixed-length dense vector). This indeed simplifies the calculation.

% Human languages, like many biological systems including families of proteins, genomes, and neurons in the brain, have significant long-range correlations that decay with a power law \citep{tagliazucchi2012criticality, mora2011biological}. However,  $n$-gram models have long-range correlations that decay exponentially which belie the essential critical behavior inherent in language \citep{pestun2017language}. Even 
% current network models like LSTM \cite{hochreiter1997long} are still hard to match long-range and higher-order statistics of natural languages \citep{lin2016critical}.

Human languages, like many biological systems, including families of proteins, genomes, and neurons in the brain, have significant long-range correlations that decay with a power law \citep{tagliazucchi2012criticality, mora2011biological}.  
Current network models like LSTMs \cite{hochreiter1997long} are hard to match long-range and higher-order statistics of natural languages \citep{lin2016critical}. 


Recently, researchers have turned to tensor network language modeling,  which contains models that exhibit correlation functions that decay with the power law \citep{pestun2017tensor,pestun2017language,miller2021tensor}. Tensor networks are, roughly, decompositions of large tensors into sets of smaller tensors and have been employed in physics, mathematics, and machine learning \citep{cohen2016expressive}. However, the so-called `tensor network language model' is either a concept that needs to be proved practically \cite{pestun2017tensor} or unsuitable in real-world language modeling tasks \citep{miller2021tensor} due to their way of modeling probabilities.  Towards making tensor network language modeling practical, we make the first step to applying it to real language modeling datasets.

As proof-of-concept work, we derive a Tensor Train Language Model (TTLM) (the simplest tensor network).  Technically, we represent a sentence based on the exponential semantic space constructed by the tensor product of word representations. The probability of the sentence is defined by the inner product of two high-dimensional tensors: the input $\Phi(X)$ and the global 
coefficients $\mathcal{A}$, and decomposed into conditional probabilities.  

Under the framework of TTLM,  we propose two variants:  TTLM-Tiny and TTLM-Large. Also, we clarify the relationship between the proposed TTLM and a series of Recurrent Neural Networks (RNNs) (i.e., Second-order RNNs \citep{goudreau1994first},  Recurrent Arithmetic Circuits (RACs) \citep{levine2018benefits}, and Multiplicative Integration RNNs (MI-RNNs) \citep{wu_multiplicative_2016}). These connections open a new eye to understanding RNNs and give some natural implementations for TTLM.


We benchmark these TTLM variants and analyze the difference in their working mechanism and behaviors. Experimental results on the language modeling task show that our TTLM variants could outperform than Vanilla-RNNs under the same training setting. These demonstrate the feasibility of TTLM.
 
 
The main contributions of our work can be summarized as follows:
\begin{itemize}
    % \item We propose a novel Tensor Train Language Model, as a first attempt to apply tensor networks on real-world language modeling tasks.
    \item We propose a novel Tensor Train Language Model, as an illustration of how tensor networks can be applied to real-world language modeling datasets. 
    \item We propose two novel TTLM variants, TTLM-Large and TTLM-Tiny, and theoretically demonstrate the relationship between  TTLM and a series of existing RNNs.
    \item Compared to Vanilla-RNNs on WikiText-2 and PTB datasets, TTLM-Large reduces perplexity by 14.3 and 16.0, respectively, and TTLM-Tiny reduces perplexity by 1.7 and 8.5, respectively.
\end{itemize}


